\chapter{Experimental Results}

This chapter presents the experimental evaluation of the proposed NLGCN
framework on four real-world network datasets. The model is compared
against three classical baseline centrality methods using Kendall Tau
correlation and Top-$N$ spread evaluation via the SIR
model \cite{ chen2012}.

\section{Datasets}

Experiments are conducted on four real-world benchmark networks of
varying sizes and domains \cite{hafiene2020, farooq2018}. Table
\ref{tab:datasets} summarizes the statistical properties of each
dataset.

\begin{table}[h]
\centering
\caption{Summary of Benchmark Datasets}
\label{tab:datasets}
\begin{tabular}{|l|c|c|c|l|}
\hline
\textbf{Dataset} & \textbf{Nodes} & \textbf{Edges} &
\textbf{Avg. Degree} & \textbf{Network Type} \\
\hline
C. elegans    & 297 & 2{,}148 & 14.46 & Biological neural network \\
Budapest      & 480 & 1{,}000 &  4.17 & Transportation network \\
US Airports   & 500 & 2{,}980 & 11.92 & Airport route network \\
CargoShipsBB  & 834 & 4{,}349 & 10.43 & Cargo shipping network \\
\hline
\end{tabular}
\end{table}

All datasets are treated as unweighted, undirected graphs. Average
degree is computed as $2 \times |E| / |V|$. The Budapest network had
11 self-loops removed during preprocessing. The datasets span a range
of network sizes (297--834 nodes) and structural characteristics,
from dense biological networks to sparse transportation graphs, enabling
a comprehensive evaluation of the model's generalization ability.

\section{Experimental Setup}

\subsection{Framework and Hardware}

All experiments are implemented in PyTorch (v2.11.0) with Python 3.10.9.
Supporting libraries include NetworkX for graph operations, NumPy and
SciPy for numerical computations, scikit-learn for preprocessing, and
Matplotlib for visualization. All experiments are run on a local CPU
(Windows) without GPU acceleration.


\subsection{Hyperparameters}

Table~\ref{tab:hyper} lists the hyperparameters used across all
experiments.

\begin{table}[h]
\centering
\caption{Hyperparameter Configuration}
\label{tab:hyper}
\begin{tabular}{|l|c|}
\hline
\textbf{Parameter} & \textbf{Value} \\
\hline
Learning Rate & 0.001 \\
Optimizer & Adam \\
Loss Function & MSELoss \\
Epochs & 300 \\
Input Channels & 6 (NLI + NGI at $L = 0, 1, 2$) \\
SIR Runs (label generation) & 500 per node \\
Infection Probability ($\beta$) & $1.5 \times \beta_c$ \\
Recovery Probability ($\mu$) & 1 \\
Input Normalization & Per-channel zero-mean, unit-variance \\
Label Normalization & Min-max to $[0,1]$, then z-score \\
\hline
\end{tabular}
\end{table}

The epidemic threshold is computed as $\beta_c = \langle k \rangle /
(\langle k^2 \rangle - \langle k \rangle)$,
and the infection probability is set to $\beta = 1.5\,\beta_c$ to
operate above the epidemic threshold, consistent with prior work
\cite{guo2026, chen2012}.

\subsection{Training Convergence}

The model converges reliably across all datasets. Table
\ref{tab:convergence} reports the MSE training loss at the start and
end of training.

\begin{table}[h]
\centering
\caption{Training Loss Convergence}
\label{tab:convergence}
\begin{tabular}{|l|c|c|}
\hline
\textbf{Dataset} & \textbf{Loss (Epoch 0)} &
\textbf{Loss (Epoch 280)} \\
\hline
C. elegans   & 1.0288 & 0.0030 \\
US Airports  & 1.8410 & 0.0018 \\
Budapest     & 0.9689 & $<$0.01 \\
CargoShipsBB & 1.0597 & $<$0.01 \\
\hline
\end{tabular}
\end{table}

The substantial reduction in loss (over 99\% reduction in all cases)
confirms that the NLGCN model successfully learns to approximate
SIR-based influence scores from structural features.

\section{Baseline Methods}

The proposed NLGCN is compared against three widely used centrality
measures \cite{mukhtar, farooq2018}:

\begin{itemize}
    \item \textbf{Degree Centrality} : Ranks
    nodes by the number of direct connections. Simple and
    computationally efficient, but ignores global network structure.

    \item \textbf{K-Shell Decomposition}: Assigns each
    node a core number via iterative pruning. Nodes in higher-$k$
    cores are considered more central. Captures some global structure
    but produces coarse rankings with many tied scores.

    \item \textbf{Betweenness Centrality} :
    Measures how often a node lies on shortest paths between other node
    pairs. Identifies bridge nodes but has $O(n^3)$ computational
    complexity and does not directly measure spreading ability.
\end{itemize}

\section{Comparison with Baseline Methods}

\subsection{Kendall Tau Correlation}

Table~\ref{tab:kendall} presents the Kendall Tau correlation
 between the rankings produced by each method and the
ground truth SIR-based rankings.

\begin{table}[h]
\centering
\caption{Kendall Tau Correlation Comparison on Benchmark Datasets}
\label{tab:kendall}
\begin{tabular}{|l|c|c|c|c|}
\hline
\textbf{Method} & \textbf{C. elegans} & \textbf{Budapest} &
\textbf{US Airports} & \textbf{CargoShipsBB} \\
\hline
Degree Centrality         & 0.7472 & 0.5537 & 0.6503 & 0.7874 \\
K-Shell                   & 0.7683 & 0.6277 & 0.6898 & 0.8246 \\
Betweenness Centrality    & 0.5723 & 0.2729 & 0.3592 & 0.5957 \\
\textbf{NLGCN} & \textbf{0.9627} & \textbf{0.9164} &
\textbf{0.8944} & \textbf{0.9237} \\
\hline
\end{tabular}
\end{table}

NLGCN achieves the highest Kendall Tau on every dataset, with scores
ranging from 0.8944 (US Airports) to 0.9627 (C. elegans). Several key
observations are noted:

\begin{itemize}
    \item K-Shell is the strongest baseline (0.6277--0.8246), followed
    by Degree Centrality (0.5537--0.7874). Betweenness Centrality
    performs worst, particularly on sparse networks (Budapest: 0.2729).

    \item NLGCN's advantage is most pronounced on the Budapest network
    (+46.0\% over K-Shell), which is the sparsest dataset (avg. degree
    4.17). This demonstrates the model's ability to capture structural
    patterns that simple centrality measures miss in low-connectivity
    settings.

    \item Table~\ref{tab:improvement} quantifies the percentage
    improvement of NLGCN over each baseline.
\end{itemize}

\begin{table}[h]
\centering
\caption{NLGCN Percentage Improvement Over Baseline Methods}
\label{tab:improvement}
\begin{tabular}{|l|c|c|c|c|}
\hline
\textbf{Baseline} & \textbf{C. elegans} & \textbf{Budapest} &
\textbf{US Airports} & \textbf{CargoShipsBB} \\
\hline
vs Degree Centrality   & +28.8\% & +65.5\% & +37.5\% & +17.3\% \\
vs K-Shell             & +25.3\% & +46.0\% & +29.7\% & +12.0\% \\
vs Betweenness         & +68.2\% & +235.8\% & +149.0\% & +55.1\% \\
\hline
\end{tabular}
\end{table}


